\chapter{งานและพลังงานกลในระบบฟิสิกส์}
\label{ch:ch07_work_energy}
\index{งานและพลังงานกล}

\begin{conceptbox}[title=วัตถุประสงค์การเรียนรู้ประจำบท]
เมื่อสิ้นสุดการศึกษาบทเรียนนี้ นักศึกษาสามารถ
\begin{enumerate}
    \item อธิบายนิยามของงานเชิงกลในทางฟิสิกส์ และคำนวณงานที่เกิดจากแรงคงที่และแรงแปรผันได้อย่างถูกต้อง
    \item อธิบายความหมายของพลังงานจลน์และประยุกต์ใช้ทฤษฎีบทงาน-พลังงาน (Work-Energy Theorem) ในการแก้โจทย์ปัญหา
    \item อธิบายความหมายของพลังงานศักย์โน้มถ่วงและพลังงานศักย์ยืดหยุ่นของสปริง พร้อมทั้งคำนวณค่าได้
    \item จำแนกความแตกต่างระหว่างแรงอนุรักษ์และแรงไม่อนุรักษ์ และวิเคราะห์ผลของแรงเสียดทานต่อพลังงานกล
    \item ใช้กฎการอนุรักษ์พลังงานกลวิเคราะห์และแก้โจทย์ปัญหาการเคลื่อนที่ในสถานการณ์ต่าง ๆ เช่น การเคลื่อนที่บนทางลาดและรางโค้ง
\end{enumerate}
\end{conceptbox}

\section{ทฤษฎีและหลักการพื้นฐาน}

\begin{bioappbox}
\textbf{ชีวพลังงานและกฎการอนุรักษ์พลังงานในเซลล์สิ่งมีชีวิต}

พลังงานในระบบชีวภาพสอดคล้องกับกฎการอนุรักษ์พลังงานอย่างเคร่งครัด โดยพลังงานแสงอาทิตย์ถูกเปลี่ยนรูปเป็นพลังงานศักย์เคมีในพันธะของน้ำตาลกลูโคสผ่านกระบวนการสังเคราะห์ด้วยแสง และถูกปลดปล่อยออกมาเป็นพลังงานกลและความร้อนในกระบวนการหายใจระดับเซลล์ (Cellular Respiration) ผ่านสารพลังงานสูง ATP การคำนวณงานทางฟิสิกส์ $W = \int \vec{F} \cdot d\vec{r}$ ยังใช้อธิบายงานในการสูบฉีดเลือดของหัวใจและการลำเลียงน้ำตาลในท่อโฟลเอ็ม (Phloem)
\end{bioappbox}

\begin{labbox}
1. ตรวจสอบตำแหน่งศูนย์ (Zero Error) ของเครื่องมือวัดทุกชนิดก่อนเริ่มบันทึกข้อมูล
2. ทำการวัดซ้ำอย่างน้อย 3 ถึง 5 ครั้งในแต่ละจุดทดลองเพื่อคำนวณหาค่าเฉลี่ยและค่าเบี่ยงเบนมาตรฐาน
3. บันทึกผลการวัดพร้อมระบุหน่วยเอสไอ (SI Units) และค่าความคลาดเคลื่อนของการวัดเสมอ
\end{labbox}

\section{ตัวอย่างโจทย์และการวิเคราะห์เชิงฟิสิกส์}

\begin{examplebox}
แรงไม่อนุรักษ์ (Non-conservative Force) คือแรงที่งานซึ่งกระทำต่อวัตถุ ขึ้นอยู่กับเส้นทางการเคลื่อนที่ ตัวอย่างที่สำคัญที่สุดคือ แรงเสียดทาน (friction) ซึ่งงานที่กระทำโดยแรงเสียดทานจะมีค่ามากขึ้นตามระยะทางที่วัตถุเคลื่อนที่ผ่าน แม้ว่าตำแหน่งเริ่มต้นและตำแหน่งสุดท้ายจะเหมือนกันก็ตาม แรงไม่อนุรักษ์มักจะทำให้พลังงานกลของระบบสูญเสียไปในรูปของพลังงานความร้อนหรือพลังงานเสียง จึงไม่สามารถนิยามพลังงานศักย์ที่สัมพันธ์กับแรงประเภทนี้ได้

การจำแนกแรงออกเป็นสองประเภทนี้มีความสำคัญอย่างยิ่งในการนำหลักการอนุรักษ์พลังงานกลไปประยุกต์ใช้ เนื่องจากกฎการอนุรักษ์พลังงานกลจะใช้ได้ก็ต่อเมื่อในระบบนั้นไม่มีแรงไม่อนุรักษ์กระทำ หรือกล่าวอีกนัยหนึ่งคือ งานที่กระทำโดยแรงไม่อนุรักษ์มีค่าเป็นศูนย์เท่านั้น

การแปลงผันระหว่างพลังงานศักย์และพลังงานจลน์บนรางไร้แรงเสียดทาน

7.5 กฎการอนุรักษ์พลังงานกล
\end{examplebox}

\section{แบบฝึกหัดทบทวนและคำถามท้ายบท}

\begin{enumerate}
    \item จงอธิบายความแตกต่างระหว่างงานที่มีค่าเป็นบวก งานที่มีค่าเป็นลบ และงานที่มีค่าเป็นศูนย์ พร้อมยกตัวอย่างสถานการณ์แต่ละกรณี
    \item คนงานคนหนึ่งดันกล่องด้วยแรง 80 นิวตัน ในแนวราบ ทำให้กล่องเคลื่อนที่ไป 6 เมตร จงหางานที่คนงานกระทำต่อกล่อง
    \item เชือกลากวัตถุด้วยแรง 120 นิวตัน ทำมุม 40$^\circ$ กับแนวราบ ทำให้วัตถุเคลื่อนที่ไป 10 เมตร จงหางานที่เกิดจากแรงลาก
    \item จงอธิบายความหมายของพลังงานจลน์ และยกตัวอย่างสถานการณ์ในชีวิตประจำวันที่แสดงถึงพลังงานจลน์อย่างชัดเจน
    \item วัตถุมวล 3 กิโลกรัม เคลื่อนที่ด้วยอัตราเร็ว 4 m/s จงหาพลังงานจลน์ของวัตถุ
    \item หากอัตราเร็วของวัตถุในข้อ 5 เพิ่มขึ้นเป็นสองเท่า พลังงานจลน์ของวัตถุจะเปลี่ยนแปลงไปอย่างไร จงอธิบายพร้อมคำนวณ
\end{enumerate}

% สิ้นสุดบทเรียน
